\begin{exercisechunk}
\item \textbf{Majority cycles in random utility models.}\suppmark
\label[exercise]{ex:l19-random-utility-cycles}
\exerciseprofile{2}{2}{2}{\exproof\quad\exconstruction\quad\exinterpretation}{%
Distinguish failure of Bradley--Terry from failure of a general random utility model.}
Let $U=(U_a,U_b,U_c)\in\R^3$ be jointly distributed, with
$\P(U_y=U_{y'})=0$ for distinct $y,y'\in\{a,b,c\}$.
The same utility law is used for every compared pair, independently of
which pair is presented. Define $p(y,y')=\P(U_y>U_{y'})$.
\begin{enumerate}[(a)]
\item Prove the triangle bounds
\[
 1\le p(a,b)+p(b,c)+p(c,a)\le2.
\]
\item Construct a distribution supported on three utility vectors for which
$p(a,b)=p(b,c)=p(c,a)=2/3$. Verify these probabilities and the absence
of ties.
\item Decide whether the probabilities $p(a,b)=p(b,c)=p(c,a)=3/4$
from \cref{sec:misspecification} admit such a random utility model.
Then construct a random utility model for the pooled probabilities
from \cref{ex:l19-pooled-evaluators}. Explain why incompatibility with
Bradley--Terry has different implications in these two examples.
You may use the identity requested in \cref{ex:l19-gumbel-comparisons}.
\end{enumerate}
\end{exercisechunk}
